\section{Recovering the lower parameter from an upper theta lift}

% BEGIN U86T1
\paragraph{U86T1. The split parameter of a nonzero pairing.}
Write $T$ for the up operator at the original localization $i$.
Suppose $\langle F,Tf\rangle\ne0$ for some $F$ of upper type $P$
and some $f$ in the lower level space of type $\sigma$. Fix any
split finite place $w$ of $E$ and put $n=r+t$. In the two general
linear groups let the components be $\sigma_w$ and $\pi_w$.
Then, in the member's local-parameter convention,
\[
 \operatorname{rec}(\pi_w)=
 \left(\operatorname{rec}(\sigma_w)\oplus
      \bigoplus_{k=0}^{t-1}\operatorname{rec}(\mu_{i,w})
                  |\cdot|^{(t-1)/2-k}\right)\otimes
                     \operatorname{rec}(\eta_w^{-1}).
 \tag{1}
\]

Realize $F$ as a definite algebraic modular form and $f$ in the
cuspidal multiplicity space with its invariant inner product.
On the upper side the member's adelic identification gives the
weight intertwiner $e_W$; its inverse identifies the weight values
of these forms. The rational group actions agree in the two standard
frames, and the member's evaluation identity transports the functions
by this same intertwiner. On the lower side the two readings give
isomorphic copies of the same finite type by the row-class
identification. The coefficient and tuple identifications are those
of M84T1 and M84T6. Thus the principal-level map in L86M1 evaluates,
up to the common nonzero metric normalization, by the global theta
integral on these forms and the prescribed Schwartz vector.
The adjoint identity $\langle F,Tf\rangle=\langle T^*F,f\rangle$
puts the lower form in the conjugate-linear variable of the
trilinear pairing (or take the complex conjugate of the entire
identity for the opposite inner-product convention).
The integral is invariant under simultaneous translation of the
kernel and forms by change of the two integration variables.
The upper quotient is compact and cuspidal rapid decay on the
lower quotient justifies this change for the theta kernel.

At $w$ the Schwartz space is spanned by local tensors with
other coordinates fixed. For each finite-adelic irreducible,
fixing the other tensor coordinates likewise gives equivariant
copies of its local component; these span the smooth vectors.
For an isotypic space take the sum of the images of the irreducible.
Successively expand the three arguments in these spans. Some triple
has nonzero value. With the other coordinates fixed the integral is
a nonzero local trilinear pairing. Translations supported at the split
place realize arbitrary matrices in both factors. In the split
Schwartz model the upper determinant contributes the scalar
$\eta_w(\det h)$; putting it on the upper vector gives
$\pi_w\otimes(\eta_w\circ\det)$. In the lower variable use its invariant Hermitian form to
identify the antilinear functional with a vector in the smooth
lower representation. The linear upper variable gives a vector in
the smooth contragredient of the twisted upper representation.
For a fixed smooth Schwartz vector its pairing is fixed by an
open compact product. Averaging each automorphic factor over its
compact subgroup factors the functional through their two spaces
of fixed vectors. Those are finite-dimensional by admissibility,
so their product dual is the tensor of their duals. Varying the
Schwartz vector gives the equivariant quotient map required by
the type-II correspondence. This gives the type-II pairing for $\sigma_w$ with splitting character
$\mu_{i,w}$.

In the explicit type-II correspondence, passing from rank $r$ to
$r+t$ inserts $t$ consecutive exponents centered at zero in the
contragredient Langlands data (A. M\'inguez, \emph{Correspondance
de Howe explicite}, Ann. Sci. ENS 41 (2008), Theorem 1).
The upper factor in the quotient is the smooth contragredient; taking its dual cancels the contragredient of the lower Langlands data.
In the rank-$(r,r+t)$ matrix model the character multiplier on
$(g,h)$ is $\mu_{i,w}(\det g)\mu_{i,w}(\det h)^{-1}$.
The determinant characters in this action replace the trivial
inserted block by $\mu_{i,w}\circ\det$: twist the lower input by
$\mu_{i,w}^{-1}$ before applying the formula and twist the inducing
factors back by $\mu_{i,w}$. The lower representation is
unitarizable: pull the positive invariant form back to a nonzero
copy of the finite type, then to the embedded local component.
The normalized product
$\sigma_w\times(\mu_{i,w}\circ\det)$ is irreducible by Bernstein's
unitary product theorem. Its parameter is the direct sum of the
parameters of its factors; that of the determinant character
consists of the displayed $t$ lines with zero monodromy. For the
upper label, H86P1 and B85U1--2 give good split supercuspidal
support at every split place: the original upper label is
spherical at a nonsplit place, and the member's reading of this
predicate includes the support condition. Twisting preserves
good support. In the irreducible product the two inducing supports
form the support of the product, so the lower support is good as
well. Thus A72F7 identifies both weak parameter readings with
their classical ones. Reading the resulting parameter and undoing the upper
twist proves (1). In particular, if the entire nonzero image
$T(\mathrm{isot}_j(\sigma))$ is in type $P$, take $F=Tf\ne0$.
This proves the original $\mathrm{adm3}$.
% END U86T1

% BEGIN U86T2
\paragraph{U86T2. The row bound and uniqueness of the upper label.}
At a split place where its parameter is unramified, the component
of a nonzero lower row is spherical by the spherical local
correspondence: the classical inverse of an unramified parameter is
spherical and has character support, on which the weak and classical
maps agree. The weak-map injectivity with one good input in A72F7
therefore identifies the lower component. It is unitary by the invariant-form argument of
U86T1, which only uses occurrence of the lower type. Write
$q=\mathrm Nw$. Each spherical unitary Satake root, and its
inverse, has absolute value at most $q^{(r-1)/2}$. Indeed, in the
spherical case of Tadi\'c's unitary dual classification a unitary
character block of length $b$ has exponents from $-(b-1)/2$ to
$(b-1)/2$. A complementary pair of such blocks adds opposite
shifts of absolute value less than $1/2$ and consumes $2b$ of the
rank. Both kinds give the required bound. Only unramified character
blocks occur in a spherical product, by the hyperspecial invariants
in normalized induction and the spherical classification. Since
$q>1$ and $r<t$, this gives the strict bound at exponent $(t-1)/2$,
also for the inverse-Satake reading. This is $\mathrm{RowsB.rowSize}$.

Suppose $P,P'$ both pair nontrivially with this lower row.
Equation (1) gives equal split parameters at every split place.
Take their base-change lists from H86P1 and the representations
of H86P2. At good split Frobenius elements these have equal
characteristic polynomials. By S86A1 the semisimple representations
are isomorphic. H86P2 now gives equal list parameters at almost
all places, including nonsplit ones, after undoing the common
unramified twist. S86A1 recovers equality of the lists up to
permutation. At every nonsplit place both labels are spherical
and have equal unramified base-change multisets by H86P1. Remove
one $1$ in odd rank. On each two-element inversion orbit the number
of pairs is its common multiplicity; at $1,-1$ the remaining
multiplicities must be halved. This determines the relative tuple
up to permutations and individual inversions, precisely the
relative Weyl action. Thus the nonsplit representations agree.
At split places use (1) and injectivity of the local correspondence
on the supported components. Taking the restricted tensor product
of these component isomorphisms, normalized on spherical vectors
almost everywhere, identifies the finite-adelic irreducibles.
The upper labels enumerate their isomorphism classes in fixed
weight, so $P=P'$ even if the representation has multiplicity.
% END U86T2

% BEGIN U86T3
\paragraph{U86T3. The whole lower image lies in a single upper type.}
At fixed $i$ the upper level space is spanned by its type spaces.
They are orthogonal in the native inner product: the member reads
this product through the definite invariant class-sum pairing.
A nonzero cross pairing between two simple copies would give a
nonzero intertwiner from one to the Hermitian dual of the other.
That dual is isomorphic to the other unitary representation, so
Schur's lemma gives orthogonality for inequivalent copies. Taking
sums and imposing the level preserves this orthogonality.

If the lower summand is nonzero, L86M4 gives $T^*Tf=\lambda f$
on it with $\lambda>0$. Choose $f\ne0$. Since
$\|Tf\|^2=\lambda\|f\|^2>0$, some upper type $P$ pairs nontrivially.
For any $f'$ in the same lower summand every component of $Tf'$
outside $P$ is zero: a nonzero orthogonal projection would provide
a second nonzero pairing, contradicting U86T2. Thus $P$ satisfies
$\mathrm{thetaOf}(i,\sigma,P)$. The nonzero image and disjointness
of the type spaces make $P$ unique. This proves
$\mathrm{RowsB.thetaOne}$. With U86T1--2 and I86A3 it gives all
four fields of $\mathrm{RowsB}$.
% END U86T3

% BEGIN U86T4
\paragraph{U86T4. The long string, the good rows and the original first conclusion.}
Fix this $P$ and a base-change list $\mathcal P$ for it.
At almost every split place all components are spherical. In (1)
there is exactly one root of absolute value $q^{(t-1)/2}$, the
extreme root from the distinguished character string. All lower
roots are smaller by U86T2, and no root of (1) is larger.
Every unramified root of a unitary cuspidal base has absolute value
strictly less than $q^{1/2}$ by S86A2. Hence a block of length
$d<t$ in $\mathcal P$ has all roots strictly smaller than
$q^{(t-1)/2}$. There must be a block of length $d\ge t$.
The total rank is $r+t<2t$, so exactly one block has this property
and its cuspidal base is a character. Unitarity makes its extreme
absolute value exactly $q^{(d-1)/2}$, whence $d=t$. Equating the
two extreme roots at every such split place and cancelling the
common positive norm factor identifies this character locally with
$\mu_i\eta^{-1}$. The rank-one comparison in S86A2 identifies it
globally, on the member's original labels.

Remove this block and twist each remaining cuspidal base by $\eta$
to form $\Sigma$. It has positive parts and rank $r$. At every
split place untwist (1) and cancel the common character string
using I86A3, direct-sum cancellation for Weil--Deligne parameters.
This gives precisely $\mathrm{HasParamW}(i,\sigma,\Sigma)$,
including ramified split places. Appending $(1,\mu_i,t)$ and
twisting by $\eta^{-1}$ recovers $\mathcal P$ up to permutation,
so all six base-change identities, including the nonsplit and
infinity identities, hold for this reconstructed list. U86T3
makes the assertion valid for every $P'$ with $\mathrm{thetaOf}$.

Finally suppose the row is good and let $\Sigma'$ be \emph{any}
list satisfying $\mathrm{HasParamW}$ for it. At each designated
place its local reading equals the lower reading. The row label is
the polynomial of the lower reading on wild invariants by the
member's finite-type comparison (M84T9 and G85M3). Goodness gives
full degree $r$, and the sum of the ranks of the parts, counted
with lengths, is $r$. Each factor must have its full positive
degree. A length $d>1$ contributes the $d$th power of that factor,
contrary to squarefreeness since norm shifts leave inertia unchanged.
If several length-one blocks remain, one has degree $a$ with
$1\le a\le r/2$. It would give a divisor over the residue field
of degree $a$ at every designated place over the first type prime.
The missing-divisor-degree clause supplies a place without such
a divisor. Thus $\Sigma'$ consists of one rank-$r$ cusp of length
one. This proof also covers lists for which conjugate self-duality
was not assumed.

Together with L86M4 and U86T3 this proves all three assertions of
$\mathrm{Conclusion42}$. H86P4 proves $\mathrm{provisionFortyOne}$
and S86A4 proves the unused middle conjunct. Hence the original
member field $\mathrm{GaloisPreconditions.Conclusion1}$ holds.
% END U86T4
